% Part: first-order-logic
% Chapter: semantics
% Section: intro-semantics

\documentclass[../../../include/open-logic-section]{subfiles}

\begin{document}

\olfileid[pt]{fol}{syn}{its}

\olsection{Introdução}

Dar o significado das expressões é o domínio da semântica. O
conceito central em semântica é o de satisfação em
!!a{structure}. !!^a{structure} dá significado aos blocos de construção
da linguagem: !!a{domain} é um conjunto não vazio de objetos. Os
quantificadores são interpretados como variando sobre esse domínio, aos !!{constant}s
são atribuídos elementos no domínio, aos !!{function}s são atribuídas
funções do !!{domain} para si mesmo, e aos !!{predicate}s são
atribuídas relações sobre o !!{domain}. O !!{domain} junto com as
atribuições ao vocabulário básico constitui !!a{structure}.
!!^{variable}s podem aparecer em !!{formula}s, e para dar uma
semântica, também temos que atribuir !!{element}s do !!{domain} a
elas---isto é uma atribuição de variáveis. A relação de satisfação,
finalmente, reúne tudo isso. !!^a{formula} pode ser satisfeita em
!!a{structure}~$\Struct{M}$ relativamente a uma atribuição de variáveis~$s$,
escrita como $\Sat{M}{!A}[s]$. Essa relação também é definida por
indução sobre a estrutura de~$!A$, usando as tabelas de verdade para os
conectivos lógicos para definir, por exemplo, a satisfação de $(!A \land !B)$ em
termos da satisfação (ou não) de $!A$ e~$!B$. Acontece então
que a atribuição de variáveis é irrelevante se !!a{formula}~$!A$
é !!a{sentence}, isto é, não tem variáveis livres, e assim podemos falar de
!!{sentence}s sendo simplesmente satisfeitas (ou não) em !!{structure}s.

Com base na relação de satisfação $\Sat{M}{!A}$ para !!{sentence}s
podemos então definir as noções semânticas básicas de validade, consequência
e satisfatibilidade. !!^a{sentence} é válida, $\Entails !A$, se toda
!!{structure} a satisfaz. Ela é consequência de um conjunto de !!{sentence}s,
$\Gamma \Entails !A$, se toda !!{structure} que satisfaz todas as
!!{sentence}s em~$\Gamma$ também satisfaz~$!A$. E um conjunto de !!{sentence}s
é satisfatível se alguma !!{structure} satisfaz todas as !!{sentence}s nele
ao mesmo tempo. Como !!{formula}s são definidas indutivamente, e
a satisfação é por sua vez definida por indução sobre a estrutura das
!!{formula}s, podemos usar indução para provar propriedades de nossa
semântica e para relacionar as noções semânticas definidas.

\end{document}
